% Copyright (C) 2018 by latexstudio <http://www.latexstudio.net>
%
% This program is free software: you can redistribute it and/or modify
% it under the terms of the GNU General Public License as published by
% the Free Software Foundation, either version 3 of the License, or
% (at your option) any later version.
%
% This program is distributed in the hope that it will be useful,
% but WITHOUT ANY WARRANTY; without even the implied warranty of
% MERCHANTABILITY or FITNESS FOR A PARTICULAR PURPOSE.  See the
% GNU General Public License for more details.
%
% You should have received a copy of the GNU General Public License
% along with this program.  If not, see <http://www.gnu.org/licenses/>.
%

\section{介绍公式的常见问题}


\faq{\cs{[} \ldots \cs{]} 与 \texttt{\$\$\ldots \$\$} 有什么区别？}{}

重定义的难度不同、造成的间距也不同。推荐使用 \cs{[} \ldots \cs{]}。 见
\url{https://www.zhihu.com/question/27589739/answer/37255728}


\faq{如何让长公式自动断行？}{}

长公式自动断行要看情况，如果是在行内模式，合理使用空格，一般可以在二元运算符处断行，如果是行间模式，推荐使用align类环境，在需要断行处添加|\\|%
%\textbackslash{} 
手动断行。


\faq{公式希腊字符如何加粗？}{}

希腊字母没有粗体，可以选择合适的数学字体。 可以使用 bm
宏包将希腊字母加粗。


\faq{极限符号下面有两个趋近该怎么写}{}

直接给出例子：

\begin{verbatim}
\documentclass{article}
\begin{document}
 \[ \lim_{n\to\infty\atop m\to\infty} \]
\end{document}
\end{verbatim}
效果如下：

\[ \lim_{n\to\infty\atop m\to\infty} \]

或者使用 \cs{substack}，代码如下：

\begin{verbatim}
\documentclass{article}
\usepackage{amsmath}
\begin{document}
\[ \lim_{\substack{n\to\infty\\ m\to\infty}} \]
\end{document}
\end{verbatim}

效果如下：

\[ \lim_{\substack{n\to\infty\\ m\to\infty}} \]
 
% \includegraphics{https://images-cdn.shimo.im/FCY4A1SeBIcwBCGT/双重极限.PNG!thumbnail}


\faq{怎样在 LaTeX 中输入引号}{}

左引号用 `（键盘1旁边那个键），右引号用 `。双引号也一样，``''。
中文条件下可以直接用中文引号（这个与编码方式和中文支持方式有关的），会有自动配对（这个和编辑器以及输入法有关的），但是如果需要用到不配对引号的情况，需要使用通用方法。


\faq{align环境默认是居中对齐吗？我在使用时，发现公式开始是居中的，后来却一直靠右断对齐，这是什么原因？}{}

%\sout{align 默认靠右对齐，所以通常加 \&
%符号，让代码左对齐。验证一下以下代码：}

\begin{verbatim}
\begin{align}
& \nabla \times H = J,\\
& \nabla \times E = - \partial _t B,\\
& \nabla \cdot B = 0.
\end{align}
\end{verbatim}

再试试把 \& 去掉什么样。

align采用的是奇偶对齐的方式，第一列右对齐，第二列左对齐，就这样右左右左依此类推，两列之间用\&分隔。


\faq{中英文标点使用规则不是很明白，尤其在公式环境里，字体和间距差别都比较大。怎样才能让正文和公式的标点统一（形状和间隔）？}{}

详见：
\url{https://link.zhihu.com/?target=http\%3A//www.moe.gov.cn/ewebeditor/uploadfile/2015/01/13/20150113092346124.pdf}

在导言区加类似命令可实现全文替换：

\begin{verbatim}
\catcode`\。=\active\newcommand{。}{. }
\end{verbatim}

或者使用 xeCJK 宏包的字符映射功能，调用 fullwidth-stop
这一映射文件，将中文空心句号映射为实心句点：

\begin{verbatim}
\documentclass{article}
\usepackage{xeCJK}
\setCJKmainfont[Mapping= fullwidth-stop]{STSong}
\begin{document}
句号。
\end{document}
\end{verbatim}


\faq{公式如何居左对齐,居右对齐？}{}

公式居左对齐在基础文档类中由 fleqn
选项控制，选择该选项后，正文公式均居左对齐，至于居右对齐，嗯，我没见过这么奇怪的格式。


\faq{公式之后解释公式符号的文字，通常是 ``符号 ------ 解释'' 这样的格式，我希望这段文字的格式是按破折号对齐，并且解释文字折行后悬挂缩进，怎样实现这样的格式？}{}

方法很多，可以列表，可以align等环境。 这里给出一个使用自定义列表的例子：

\begin{verbatim}
\usepackage{ifthen}
\newcounter{qlst}
\newenvironment{EqDesc}[2][式中]{%
\begin{list}{}
 {%
  \usecounter{qlst}
  \settowidth{\labelwidth}{#1，\ \ #2\ --- \ }
  \setlength{\labelsep}{0pt}
  \setlength{\leftmargin}{\labelwidth}
  \setlength{\rightmargin}{0em}
  \setlength{\parsep}{0ex}
  \setlength{\itemsep}{0ex}
  \setlength{\itemindent}{0em}
  \setlength{\listparindent}{0em}
  \renewcommand{\makelabel}[1]
    {\stepcounter{qlst}
     \ifthenelse{\value{qlst}>1}{\hfill ##1\ --- \ }{#1，\hfill ##1\ --- \ }
    }
 }%
}%
{\end{list}}
\end{verbatim}

EqDesc
环境有两个参数，第一个为可选参数，是解释公式符号前的引导词，默认是``式中''，第二个参数是样本符号，可以选择一个列表中宽度最大的符号。条目 \cs{item} 有一个可选参数（实际使用是必选参数），内容是要说明的符号。使用如下：

\begin{verbatim}
\[ a^2+b^2=c^2 \]
\begin{EqDesc}[其中]{$a$}
   \item[$a$] 三角形的一条直角边；
   \item[$b$] 三角形的另一条直角边；
   \item[$c$] 三角形的斜边。
\end{EqDesc}
\end{verbatim}

{\centerline{\includegraphics[width=0.8\linewidth]{include/images/hang}}}


\faq{行内公式的情况下如何让sum prod这些运算符的上下标在头上和脚下？}{}

这样处理行内公式的上下标会导致段落行距不整齐，不符合 \LaTeX{}
的审美。如果彻底放弃审美，可以使用 \cs{limits} 命令，如：

\begin{verbatim}
$\sum\limits_{i=1}^n \quad \prod\limits_\epsilon$
\end{verbatim}

效果是这样的：$\sum\limits_{i=1}^n \quad \prod\limits_\epsilon$

\faq{如何将积分的上限标放在积分号的上下两侧？}{}

积分号的上下限放置在积分号的右侧是英美国家和 \LaTeX{}
的排版习惯，通常无需处理。
如果你很确定需要按照 ISO 80000-2:2009 或者 GB 3102.11-93 的规定排版积分号，可以：

\begin{enumerate}
\def\labelenumi{\arabic{enumi}.}

\item
  在调用 amsmath 宏包时添加 intlimits 选项；
\item
  \texttt{\cs{def}\cs{int}\{\cs{intop}\}}
\item
  如果使用 unicode-math 宏包，
\end{enumerate}

\begin{verbatim}
\removenolimits{%
  \int\iint\iiint\iiiint\oint\oiint\oiiint
  \intclockwise\varointclockwise\ointctrclockwise\sumint
  \intbar\intBar\fint\cirfnint\awint\rppolint
  \scpolint\npolint\pointint\sqint\intlarhk\intx
  \intcap\intcup\upint\lowint
}
\end{verbatim}


\faq{如何自定义数学运算符，然后让下标放在脚下？}{}

借助 amsmath 包（实际上是amsmath自动调用的amsopn宏包）的
\cs{DeclareMathOperator*} 命令，可以在导言区定义需要的数学运算符（需要注意加不加*是有区别的）。例如

\begin{verbatim}
\DeclareMathOperator*{\esssup}{ess\,sup}
\end{verbatim}

\[
\esssup_{x\in {R}}
\]

对于只是偶尔用到的运算符，也可以直接在数学模式中使用新算符定义命令。

例如
|\operatorname*{ess\,sup}_{x\in {R}}|

\[
\operatorname*{ess\,sup}_{x\in {R}}
\]

\faq{在数学公式中，编辑等式时，每一行需要等号和等号对其，这时使用了\textbackslash{}begin\{displaymath\} \textbackslash{}begin\{split\}环境，但是呢，这些整体都是居中的，我想让式子靠左，怎么实现呢？}{}
这种情况可以考虑在文档类选项中使用fleqn选项。需要对齐的多行公式，推荐使用amsmath宏包的align等环境。

\faq{行列式变换过程中，我们一般是在中间的箭头上表示出变换的方式，如何才能在长箭头上打出多行内容？}{}
可以看看amsmath宏包提供的长箭头命令。

|\xleftarrow[下方公式]{上方公式}|   |\xrightarrow[下方公式]{上方公式}|。

\faq{如何输出反斜杠？}{}
方式很多，例如

\begin{verbatim}
    \textbackslash 或者 \verb|\|
\end{verbatim}


\faq{对equation环境下的公式、图片编号按章节、小节进行重新定义}{}

公式编号使用amsmath宏包的numberwithin设定关联即可。图片编号可以使用caption宏包自定义格式。

\faq{带括号的公式有哪些折行方式？在括号中间折行。}{}
使用| \left( \right) |括起来的内容不能直接折行，但是可以变通，比如使用| \left( \right. | 和 | \left. \right) |，
这时候需要借助幻影、支柱等方法处理各行内容高度不同带来的括号大小不同问题。

\faq{怎么在极限符号下面打出x→0 }{}
很简单的基本指令就行，例如|\lim\limits_{x \to 0}|得到$\lim\limits_{x \to 0}$

\faq{\LaTeX{}敲的公式可以插入word吗？}{}
可以，但是不推荐。
